% 【基础正文】所属：第7章（由 chapters/revised/07-value.tex \input 引入）
\minitopic{价值问题的三个层次}第一层问知识为何比假信念有价值，答案相对容易：真信念通常更好地指导行动。第二层问知识为何比单纯真信念更有价值。第三层问知识是否具有超过各组成部分价值总和的独特价值。可靠主义若只说可靠过程更可能产生真信念，就必须解释：在某个信念已经为真的情况下，可靠来源还增加什么。 这被称为\term{沼泽化问题}{swamping problem}：若真理是唯一最终认识价值，真值一旦到手，过程价值似乎被淹没\parencite{kvanvig2003}。回应可把可靠性视为工具价值、能力表现视为成就价值，或否认真理是唯一终极价值。

\minitopic{成就与归功}德性认识论认为，知识是因认知能力而取得的成功。成功与能力连接形成\term{成就}{achievement}，像靠技艺完成攀登比被风吹上山顶更值得赞赏。这解释了知识何以优于猜中。 但知识不总是高度成就。询问路人得知车站方向很容易，也似乎是知识；一个能力很强的人在恶劣环境中合理却错误，也值得认识赞许。于是价值可能分布于真、理性、能力和责任多个维度，而不是全部压在知识标签上。

\minitopic{理解不同于知道很多事实}\term{理解}{understanding}通常涉及把握各要素之间的解释、因果或依赖关系。一个人可以背诵气候事实，却不理解系统为何如此变化；也可能借助略有理想化的模型获得理解，即使模型中的个别命题并非严格真实\parencite{grimm2006,riggs2003}。 理解是否是一种知识？若“理解为什么$p$”包含知道解释，就有紧密联系；但理解具有整体性、程度性和认知把握要求。它似乎比孤立命题知识更能容忍近似，却更要求主体能回答反事实问题、重新组织材料和作出解释。

\minitopic{智慧与认识目标多元性}智慧不仅是拥有更多真信念，还涉及判断何者重要、承认限制、协调价值并在情境中行动。一个数据库可以储存大量事实，却不因此智慧。认识生活可能追求真理、避免错误、理解、好奇满足、自主、公正与智慧；这些目标有时冲突。过度避免错误会使人什么都不信，过度追求新奇会增加轻信。 因此，认识规范需要权衡而非单一最大化。教育若只奖励答对，会鼓励猜测和短期记忆；若重视解释与错误修正，才能培养持续能力。

\begin{comparison}
儒家“智”与“知”常处于伦理实践和关系判断中，庄子则警惕以有限区分穷尽世界\parencite{analects2003,zhuangzi2009}。它们提示我们区分信息占有与恰当把握。但若把智慧直接等同于当代“理解”理论，会丢失修身和生命取向。较好的比较问题是：何种认知成就值得追求，价值由真理、能力还是生活实践赋予？
\end{comparison}

\minitopic{真理价值是否足够}如果每个真信念都有同等价值，记住一串无关尘埃数目似乎与理解疾病原因一样重要。价值真理一元论可以回答：真理本身都有基本价值，但相关性、行动需要和系统联系赋予额外工具价值。多元论则把理解、好奇、自主和公正视为不可完全还原的认识价值。 多元价值会冲突。为了快速行动，我们可能接受较粗模型；为了深刻理解，则容忍暂时不确定。公共数据库最大化信息可得，却可能伤害隐私。认识论评价因此需要说明目标和约束，不能只说“更多知识总是更好”。

\minitopic{理解的测试}主体若理解一个系统，通常能解释为何发生、预测干预后变化、识别哪些因素无关，并把解释迁移到新案例。复述教科书句子不足以证明理解。反事实提问特别有用：“若氧气减少，会怎样？”“若证据不再独立，结论如何变？” 理解具有程度。新手掌握简化模型，专家掌握适用边界和异常。理想化不必摧毁理解，只要主体知道哪些方面被省略、模型在目标问题上保留何种依赖。把近似模型当作字面完整真相才会误导。

\minitopic{案例工作坊：会做题但不理解}学生记住贝叶斯公式并代入正确，却无法解释为何基础率进入分母；另一学生算术出错，但能画出自然频率树并指出条件概率方向。前者拥有某种程序能力和正确答案，后者拥有部分结构理解。教学评价若只看最终数字，会错配认识价值。 较好的评分把答案、推理、解释和错误诊断分开。允许学生修正也有认识意义：一次答对可能幸运，能够定位错误并迁移到新题更能显示稳定能力。

\minitopic{价值问题需要最小对照而非抽象赞美}要判断知识比真信念多出什么，应构造两位主体对同一命题都相信且为真，只让一位具有可靠根据、能力成就或理解。随后分别比较他们抵抗反证、迁移到新问题、获得他人信任与合理行动的表现。若所有优势都来自新增事实而非认识方式，对照就失败；若差异只在道德赞许，也尚未说明认识价值。最小对照迫使理论指出附加价值的具体位置与可观察后果。
